\documentclass[11pt,a4paper]{article}
\usepackage[utf8]{inputenc}
\usepackage[margin=1in]{geometry}
\usepackage{hyperref}
\usepackage{booktabs}
\usepackage{listings}
\usepackage{xcolor}
\usepackage{fancyhdr}
\usepackage{titlesec}
\usepackage{enumitem}

% Setup links
\hypersetup{
    colorlinks=true,
    linkcolor=blue,
    filecolor=magenta,      
    urlcolor=cyan,
    pdftitle={MarketPulse System Architecture Documentation},
    pdfpagemode=FullScreen,
}

% Header and Footer styling
\pagestyle{fancy}
\fancyhf{}
\rhead{MarketPulse Pipeline Documentation}
\lfoot{\copyright\ 2026 Alphabridge}
\rfoot{Page \thepage}

% Section styling
\titleformat{\section}
  {\normalfont\Large\bfseries}{\thesection}{1em}{}[{\titlerule[0.8pt]}]

\titleformat{\subsection}
  {\normalfont\large\bfseries}{\thesubsection}{1em}{}

% Listings styling for DDL or config snippets
\definecolor{codegreen}{rgb}{0,0.6,0}
\definecolor{codegray}{rgb}{0.5,0.5,0.5}
\definecolor{codepurple}{rgb}{0.58,0,0.82}
\definecolor{backcolour}{rgb}{0.95,0.95,0.92}

\lstdefinestyle{mystyle}{
    backgroundcolor=\color{backcolour},   
    commentstyle=\color{codegreen},
    keywordstyle=\color{magenta},
    numberstyle=\tiny\color{codegray},
    stringstyle=\color{codepurple},
    basicstyle=\ttfamily\footnotesize,
    breakatwhitespace=false,         
    breaklines=true,                 
    captionpos=b,                    
    keepspaces=true,                 
    numbers=left,                    
    numbersep=5pt,                  
    showspaces=false,                
    showstringspaces=false,
    showtabs=false,                  
    tabsize=2
}
\lstset{style=mystyle}

\begin{document}

% ─── TITLE PAGE ──────────────────────────────────────────────────────────────
\begin{titlepage}
    \centering
    \vspace*{3cm}
    {\Huge \bfseries MarketPulse Pipeline \\ [0.4cm] System Architecture Documentation} \\
    \vspace{2cm}
    {\Large \bfseries Husnain Ahmed} \\
    {\large Data Engineering Intern} \\
    {\large \bfseries ALPHABRIDGE} \\
    \vspace{2cm}
    {\large July 17, 2026} \\
    \vfill
    \begin{minipage}{0.8\textwidth}
        \centering
        \itshape
        A comprehensive overview of the end-to-end data engineering, machine
        learning, and production serving architecture for the MarketPulse financial
        platform, implemented using AWS and Databricks.
    \end{minipage}
    \vspace*{2cm}
\end{titlepage}

\newpage
\tableofcontents
\newpage

% ─── SECTION 1: PROJECT OVERVIEW ──────────────────────────────────────────────
\section{Project Overview}
MarketPulse is a production-grade, end-to-end global markets data platform designed to ingest, process, and serve financial market data at scale. The platform architecture is designed to support both high-frequency streaming market feeds and large-scale historical batch analytics using a modern medallion lakehouse paradigm. 

The system is built on a robust hybrid cloud architecture utilizing AWS serverless storage and streaming components alongside Databricks cluster computing. The pipeline ingests multi-modal data (stock and cryptocurrency trades, macroeconomic indices, and financial headlines), cleanses and transforms it through structured lakehouse layers (Bronze $\rightarrow$ Silver $\rightarrow$ Gold), runs daily sentiment analysis using FinBERT, trains return-forecasting machine learning models, and serves outputs via real-time caches, data warehouses, and interactive Streamlit applications.

% ─── SECTION 2: DATA EXTRACTION ZONE ──────────────────────────────────────────
\section{Data Extraction Zone}
This layer handles the complex ingestion of both high-frequency streaming trade data and historical batch data from disparate financial, economic, and news APIs. It acts as the primary data gateway.

\subsection{External Data Sources}
To ensure the downstream analytics and machine learning engines receive a multi-dimensional view of the market, data is collected from four distinct sources:
\begin{itemize}[leftmargin=*]
    \item \textbf{Alpaca API (Stocks \& Crypto):} Serves as the primary source for streaming market quotes and trade events. Alpaca WebSockets push high-frequency trade ticks into the streaming queue.
    \item \textbf{Yahoo Finance / yfinance API:} Acts as the historical data bootstrap. It is polled during execution to download a 5-year retrospective of daily OHLCV (Open, High, Low, Close, Volume) bars.
    \item \textbf{FRED (Federal Reserve Economic Data) API:} Polled periodically to download key macroeconomic indicators (e.g., inflation indices, interest rates, GDP growth) to provide macro context to models.
    \item \textbf{NewsAPI:} Financial news headlines are pulled continuously and processed via FinBERT, a domain-adapted NLP model, to classify headlines into numeric sentiment scores.
\end{itemize}

\subsection{AWS Streaming \& Storage Integration}
Data is staged in the AWS cloud ecosystem prior to processing:
\begin{itemize}[leftmargin=*]
    \item \textbf{AWS Kinesis Data Streams:} A serverless, auto-scaling streaming queue configured in On-Demand mode. It acts as a buffer for the live quote stream from the Alpaca WebSocket producers.
    \item \textbf{Amazon S3 Data Lake (ADLS Gen2 equivalent):} Served as the central enterprise data lake. The bucket contains three folder partitions representing the lakehouse stages: \texttt{s3://marketpulse-datalake/bronze/}, \texttt{s3://marketpulse-datalake/silver/}, and \texttt{s3://marketpulse-datalake/gold/}.
\end{itemize}

% ─── SECTION 3: ARCHITECTURAL DEVIATIONS & COST OPTIMIZATION ──────────────────
\section{Architectural Deviations \& Cost Optimization}
The architecture departs from standard Supervisor guidelines to optimize cloud costs, lower serving latencies, and increase modularity.

\subsection{Serving Layer: Streamlit App vs. AWS QuickSight}
While the initial project specification proposed using AWS QuickSight connected directly to Amazon Redshift, the implementation provides a dual-serving interface:
\begin{itemize}[leftmargin=*]
    \item \textbf{Streamlit Dashboard:} A custom Python dashboard connected to Databricks SQL Warehouse and DynamoDB. It provides real-time tick monitoring and backtest reports with no user-licensing fees.
    \item \textbf{QuickSight Integration:} Redshift Serverless tables are made available for drag-and-drop corporate reporting. By caching active dashboard queries in the serverless Streamlit app, we minimize active Redshift warehouse hours, saving up to 80\% in serving costs.
\end{itemize}

\subsection{MLOps Layer: Decoupled MLflow Registry vs. AWS SageMaker}
Rather than utilizing fully managed AWS SageMaker pipelines:
\begin{itemize}[leftmargin=*]
    \item Ephemeral python containers run training logic (\texttt{train.py}), logging parameters and metrics directly into a local, decoupled \textbf{MLflow Registry} (\texttt{mlruns.db}).
    \item Storing the verified model parameters in an open-source MLflow database hosted on a basic EC2 server eliminates SageMaker continuous endpoint hosting fees.
\end{itemize}

% ─── SECTION 4: DATABRICKS DATA TRANSFORMATION ZONE ───────────────────────────
\section{Databricks Data Transformation Zone}
The processing layer enforces the Medallion lakehouse pattern using Databricks cluster computing:
\begin{itemize}[leftmargin=*]
    \item \textbf{S3 Bronze Layer (Raw):} Initial delta tables where streaming ticks and batch files are written as-is, with basic deduplication based on unique timestamp and symbol.
    \item \textbf{S3 Silver Layer (Cleaned):} Spark jobs read Bronze tables, clean data types, join with reference sector files, handle null fields, and run news sentiment inference.
    \item \textbf{S3 Gold Layer (Analytics-Ready):} Computes multi-period technical indicators (RSI, MACD, Bollinger Bands, Moving Averages) and merges them with macro indices.
    \item \textbf{dbt Core (Databricks SQL):} Governs the transformation. Instead of ad-hoc SQL, dbt runs staging and marts models on the Databricks SQL Warehouse, running data quality tests to validate results.
\end{itemize}

% ─── SECTION 5: MLOPS \& TRAINING ZONE ────────────────────────────────────────
\section{MLOps \& Training Zone}
This layer isolates the heavy computational training workloads:
\begin{itemize}[leftmargin=*]
    \item \textbf{Training Host (train.py):} Local Docker containers query historical features from the Gold Delta tables, train an LSTM network to forecast returns, and evaluate performance.
    \item \textbf{MLflow Registry:} Logs hyperparameter runs, evaluation curves, and packages the champion model.
\end{itemize}

% ─── SECTION 6: PRODUCTION SERVING ZONE ───────────────────────────────────────
\section{Production Serving Zone}
Handles live serving and client requests with sub-millisecond latency:
\begin{itemize}[leftmargin=*]
    \item \textbf{Amazon DynamoDB Table:} Acts as the low-latency real-time feature store cache. An active Kinesis consumer parses incoming quotes and updates DynamoDB.
    \item \textbf{Amazon Redshift Serverless:} Acts as the analytical data warehouse. Historical gold marts are loaded into Redshift via \texttt{AWS\_sql\_loader.py}.
\end{itemize}

% ─── SECTION 7: ORCHESTRATION \& AUTOMATION ZONE ──────────────────────────────
\section{Orchestration \& Automation Zone}
A local Apache Airflow container host manages and schedules the pipelines:
\begin{itemize}[leftmargin=*]
    \item \textbf{DAG 1 - Ingest (1\_ingest\_bronze):} Triggers batch python producers to download data and upload raw files to S3 Bronze.
    \item \textbf{DAG 2 - Process (2\_process\_silver\_gold):} Coordinates PySpark scripts to clean and build Silver/Gold delta tables.
    \item \textbf{DAG 3 - Transform (3\_transform\_dbt):} Executes \texttt{dbt run} and \texttt{dbt test} on Databricks SQL.
    \item \textbf{DAG 4 - Quality (4\_data\_quality\_checks):} Triggers Great Expectations data verification.
    \item \textbf{DAG 5 - ML (5\_ml\_nightly\_retrain):} retrains model, logs runs, and updates Redshift warehouse.
\end{itemize}

% ─── SECTION 8: SYSTEM SYNTHESIS ──────────────────────────────────────────────
\section{Comprehensive System Synthesis \& Architectural Mechanics}

\subsection{End-to-End Operational Workflow \& Communication Protocols}
The flow of data and execution mechanics across the platform is structured as follows:
\begin{enumerate}
    \item \textbf{Streaming Ingest:} WebSocket clients connect to Alpaca, sending quotes via JSON over TCP. A Python Kinesis producer publishes records to Kinesis.
    \item \textbf{Batch Ingest:} Airflow triggers \texttt{ingest\_*.py} scripts to download daily historical OHLCV data from yfinance, FRED, and NewsAPI via HTTPS, writing Parquet data to the S3 Bronze layer.
    \item \textbf{Durable Buffering:} Real-time trade quotes are asynchronously buffered inside AWS Kinesis Stream. Concurrently, a Kinesis consumer parses the stream and updates the DynamoDB table.
    \item \textbf{Databricks Processing:} Databricks Spark clusters read raw delta tables from S3 Bronze, clean and validate schemas, and write standardized outputs to S3 Silver. Technical indicators are computed and written to S3 Gold.
    \item \textbf{dbt Transformations:} dbt compiles and runs SQL models on the Databricks SQL Warehouse, outputting gold feature marts.
    \item \textbf{MLflow Tracking:} A container task pulls gold tables, trains the LSTM forecaster, and registers the model to MLflow.
    \item \textbf{Redshift Sync:} Gold tables are copied from S3 to Amazon Redshift Serverless via the Redshift Data API.
    \item \textbf{BI Consumption:} Streamlit dashboard queries Databricks SQL (historical charts) and DynamoDB (live quote ticker), while QuickSight queries Redshift.
\end{enumerate}

\subsection{Granular Tool Functional Objectives}
\begin{table}[htbp]
\centering
\caption{System Component Matrix}
\begin{tabular}{lll}
\toprule
\textbf{Technology} & \textbf{Role in Pipeline} & \textbf{Objective} \\
\midrule
Alpaca / yfinance & Data Source Ingestion & Multi-modal financial telemetry feed \\
AWS Kinesis Streams & Streaming Buffer & Real-time tick ingestion queue \\
Amazon S3 & Medallion Data Lake & Decoupled storage for Bronze/Silver/Gold layers \\
Databricks & Cluster Computing & Big data Spark processing \& dbt Core SQL execution \\
dbt Core & SQL Transformation & Version-controlled semantic indicators (RSI, MACD) \\
Apache Airflow & Pipeline Orchestration & Automated DAG scheduling \& error notification \\
MLflow Registry & MLOps Governance & Experiment tracking \& registered champion model catalog \\
Amazon DynamoDB & Low-Latency Cache & Under 10ms real-time feature query store \\
Amazon Redshift & Analytical Data Warehouse & Serving engine for QuickSight business intelligence \\
Streamlit App & Interactive Front-end & Real-time quote feed \& backtest visualization UI \\
\bottomrule
\end{tabular}
\end{table}

\end{document}
